% ---------------------------------------------------------------------------
% eq_style.tex — HERO-EQUATION TYPOGRAPHY for the workshop deck.
%
% Public interface
%   \heroeq[TAG]{<math>}    full-width display at \large: thin Shunt accent
%                           rule on the LEFT edge only (no full box),
%                           >= 0.35cm vertical air above and below, optional
%                           small Mute eyebrow TAG (THEOREM / BOUND / MODEL).
%   \heroeq*[TAG]{<math>}   the same at \Large — reserve for key theorems.
%   \term[<color>]{<expr>}{<label>}
%                           semantic term annotation: colors <expr> with a
%                           palette color and names its role beneath via an
%                           underbrace, label in \scriptsize Mute sans.
%
% Usage rules
%   * <math> is bare math material (no \[..\]); multi-line displays go
%     through \begin{aligned}...\end{aligned} inside the argument.
%   * \term belongs at the TOP display level, never inside \frac or a
%     script: the label is centered under its brace at natural width, so
%     adjacent annotated terms can never collide — the wider of expression
%     and label sets the slot.
%   * When two adjacent expressions differ in depth (descenders,
%     subscripts), equalize their brace heights with \vphantom, e.g.
%     \term[Shunt]{e_i\vphantom{\delta_n}}{local eligibility}.
%   * Spacing conventions: \, before differentials, \! to tighten crowded
%     subscripts; no inline math taller than the line in body text.
%
% Load AFTER credit_tree_lib.tex (palette, etoolbox, tikz already active).
% Replaces the boxed \eqpanel look on content slides; \eqpanel remains
% available for backup frames.
% ---------------------------------------------------------------------------

\newsavebox{\heroeqbox}

\makeatletter
\newcommand{\heroeq}{\@ifstar{\heroeq@size{\Large}}{\heroeq@size{\large}}}
\newcommand{\heroeq@size}[1]{%
  \@ifnextchar[{\heroeq@body{#1}}{\heroeq@body{#1}[]}}
\long\def\heroeq@body#1[#2]#3{%
  \par\vspace{0.35cm}\noindent
  \begingroup
  \sbox{\heroeqbox}{%
    \begin{minipage}{\dimexpr\linewidth-0.62cm\relax}
      \ifstrempty{#2}{}{%
        {\fontsize{6.5}{8}\selectfont\bfseries\color{Mute}%
         \MakeUppercase{#2}\par}%
        \vspace{0.17cm}}%
      \centering
      % One unbreakable box: an equation wider than the panel surfaces as an
      % Overfull \hbox that the build gate rejects, instead of silently
      % wrapping after a relation sign.
      \mbox{#1\(\displaystyle #3\)}\par
    \end{minipage}}%
  {\color{Shunt}\rule[-\dp\heroeqbox]{1.5pt}%
    {\dimexpr\ht\heroeqbox+\dp\heroeqbox\relax}}%
  \hspace{0.40cm}%
  \usebox{\heroeqbox}%
  \endgroup
  \par\vspace{0.35cm}}
\makeatother

% \term[<color>]{<expr>}{<label>} — the brace is drawn in soft Mute so the
% colored expression stays the loudest element; the role label sits beneath
% in \scriptsize Mute sans (fixed size, independent of math style).
% The label carries 0.3em of padding on each side, so two adjacent
% annotated terms always keep visible air between their labels.
\newcommand{\term}[3][Ink]{%
  {\color{Mute!70}%
   \underbrace{\textcolor{#1}{#2}}_{%
     \textcolor{Mute}{\text{\scriptsize\hspace{0.3em}#3\hspace{0.3em}}}}}}
